% Full title: Essential Spectrum for Two Different Periodic Ends
% Purpose:
%   Prove the essential-spectrum formula for a scalar strip with two different
%   periodic ends and supply every localization lemma used by the proof.
% Main workflow:
%   Establish local compactness and cutoff commutator bounds, derive a Zhislin
%   criterion, construct one-sided periodic Weyl sequences, and prove both
%   inclusions in the essential-spectrum identity.
% Based on:
%   The Zhislin-sequence mechanism of De Nittis et al. (2021) and the periodic
%   Floquet--Bloch characterization of Fliss (2013).
% Main changes:
%   Adapts the one-dimensional first-order argument to the present
%   two-dimensional second-order weighted strip operator with exact periodic
%   tails and possibly different transverse periods.
% Mathematical goal:
%   Prove Theorem~\ref{thm:two-ended-essential-spectrum} without assuming the
%   essential-spectrum union formula.

\section{Essential spectrum for two different periodic ends}
\label{app:essential-spectrum}

Throughout this appendix, fix
$\beta\in[-\pi/d,\pi/d)$ and abbreviate
\[
  A:=A(\beta),\qquad A^\star:=A^\star(\beta),
  \qquad \star\in\{-,+\}.
\]
The operators act respectively in
$H=L^2(B,\rho\,\dd x\dd y)$ and
$H^\star=L^2(B,\rho^\star\,\dd x\dd y)$.  Define
\begin{equation}
  H_0:=H^-\oplus H^+,
  \qquad A_0:=A^-\oplus A^+.
  \label{eq:appendix-free-operator}
\end{equation}
Let $Q$ denote multiplication by $x$ on a strip space and let
$Q_0:=Q\oplus Q$ on $H_0$.  Choose $R_0>0$ so large that
\[
  \overline{\mathcal C^0}\subset(-R_0,R_0)\times[0,d],
  \qquad
  \rho=\rho^-\text{ for }x<-R_0,
  \qquad
  \rho=\rho^+\text{ for }x>R_0.
\]
The last two equalities hold by construction of the two periodic half-guides.
All norms below may be compared with their unweighted counterparts because of
the common positive upper and lower bounds on the coefficients.
For a scalar operator $T$, write $\rho_T=\rho$ when $T=A$ and
$\rho_T=\rho^\star$ when $T=A^\star$; formulas for $A_0$ are understood
componentwise.

\begin{lemma}[Local compactness and cutoff commutators]
\label{lem:appendix-local-compactness}
Let $T$ be any one of $A,A^-,A^+$, or $A_0$, acting in its corresponding
Hilbert space $H_T$, and let $Q_T$ be the corresponding position operator.
Then:
\begin{enumerate}[label=\textup{(\roman*)},leftmargin=*]
  \item for every bounded interval $K\subset\R$,
  \[
    \mathbf 1_K(Q_T)(T+\ii)^{-1}:H_T\longrightarrow H_T
    \quad\text{is compact};
  \]
  \item if $\zeta\in C_c^\infty(\R;[0,1])$ satisfies
  $\zeta=1$ on $[-1,1]$ and $\zeta=0$ outside $[-2,2]$, and
  $\zeta_R(x):=\zeta(x/R)$, then
  \begin{equation}
    \bigl\|[T,\zeta_R(Q_T)](T+\ii)^{-1}\bigr\|_{\mathcal B(H_T)}
    \longrightarrow0
    \qquad(R\to\infty).
    \label{eq:appendix-commutator-limit}
  \end{equation}
\end{enumerate}
\end{lemma}
\begin{proof}
We first treat a scalar component.  Its closed form is
\[
  a_T(u,v)=\int_B\nabla u\cdot\nabla\overline v\,\dd x\dd y,
  \qquad D(a_T)=H^1_\beta(B).
\]
If $u=(T+\ii)^{-1}f$, nonnegativity and the variational identity give
\[
  \|u\|_{H_T}\leq\|f\|_{H_T},
  \qquad
  \|\nabla u\|_{L^2(B)}^2
  =\operatorname{Re}\langle f,u\rangle_{H_T}
  \leq\|f\|_{H_T}^2.
\]
Thus the resolvent maps boundedly into $H^1_\beta(B)$.  Restriction to
$K\times(0,d)$ followed by the Rellich embedding
$H^1(K\times(0,d))\hookrightarrow L^2(K\times(0,d))$ is compact.  The
coefficient bounds make the weighted and unweighted local $L^2$ norms
equivalent, proving (i).  The direct-sum case follows componentwise.

Multiplication by a smooth function of $x$ preserves the operator domain:
for $u\in H^1(\Delta,B)$,
\[
  \Delta(\eta u)=\eta\Delta u+2\eta'\partial_xu+\eta''u\in L^2(B),
\]
and the horizontal $\beta$-quasiperiodic trace conditions are unchanged.
Consequently, on $D(T)$,
\begin{equation}
  [T,\zeta_R]u
  =-\frac1{\rho_T}
    \left(2\zeta_R'\partial_xu+\zeta_R''u\right),
  \label{eq:appendix-commutator-formula}
\end{equation}
with the formula interpreted componentwise for $A_0$.  Since
$\|\zeta_R'\|_\infty=O(R^{-1})$ and
$\|\zeta_R''\|_\infty=O(R^{-2})$, the preceding resolvent estimate yields
\[
  \bigl\|[T,\zeta_R](T+\ii)^{-1}f\bigr\|_{H_T}
  \leq C\left(R^{-1}+R^{-2}\right)\|f\|_{H_T}.
\]
This proves (ii).
\end{proof}

\begin{lemma}[Zhislin criterion on the strip]
\label{lem:appendix-Zhislin}
For any operator $T$ in
\Cref{lem:appendix-local-compactness} and any $\lambda\in\R$, the following
are equivalent:
\begin{enumerate}[label=\textup{(\roman*)},leftmargin=*]
  \item $\lambda\in\sigma_{\mathrm{ess}}(T)$;
  \item there exists $u_m\in D(T)$ such that
  \begin{equation}
    \|u_m\|_{H_T}=1,\qquad
    \mathbf 1_{[-m,m]}(Q_T)u_m=0,\qquad
    \|(T-\lambda)u_m\|_{H_T}\longrightarrow0.
    \label{eq:appendix-Zhislin-sequence}
  \end{equation}
\end{enumerate}
For $T=A_0$, the support condition is imposed on both components.
\end{lemma}
\begin{proof}
A sequence satisfying (ii) converges weakly to zero: this is immediate first
against compactly supported $L^2$ functions and then against arbitrary
functions by density.  It is therefore a singular Weyl sequence, which proves
(i).

Conversely, let $v_j$ be a singular Weyl sequence for $T$ and $\lambda$.
For every fixed $R$, local compactness gives
\[
  \zeta_R(Q_T)v_j
  =\zeta_R(Q_T)(T+\ii)^{-1}(T+\ii)v_j
  \longrightarrow0
  \quad\text{strongly as }j\to\infty.
\]
Indeed,
$(T+\ii)v_j=(T-\lambda)v_j+(\lambda+\ii)v_j$ converges weakly to zero and
is bounded.  Choose a subsequence, still denoted $v_m$, such that
\[
  \|\zeta_m(Q_T)v_m\|_{H_T}\leq m^{-1},\qquad
  \|(T-\lambda)v_m\|_{H_T}\leq m^{-1}.
\]
Set $w_m=(1-\zeta_m(Q_T))v_m$.  Then $w_m$ vanishes for $|x|\leq m$ and
$\|w_m\|_{H_T}\to1$.  Moreover,
\[
  (T-\lambda)w_m
  =(1-\zeta_m)(T-\lambda)v_m-[T,\zeta_m]v_m.
\]
The first term tends to zero.  For the second, write
\[
  [T,\zeta_m]v_m
  =[T,\zeta_m](T+\ii)^{-1}(T+\ii)v_m.
\]
The first factor tends to zero by
\eqref{eq:appendix-commutator-limit}, while the second is bounded.  Normalizing
$w_m$ gives \eqref{eq:appendix-Zhislin-sequence}.  This is the strip analogue
of the Zhislin construction in
\cite[Lemmas~3.11--3.12]{DeNittisEtAl2021}.
\end{proof}

\begin{lemma}[Periodic Weyl sequences placed in a prescribed end]
\label{lem:appendix-one-sided-Weyl}
Let $\star\in\{-,+\}$ and
$\lambda\in\sigma(A^\star)$.  There are normalized
$v_m^\star\in D(A^\star)$ such that
\[
  \|(A^\star-\lambda)v_m^\star\|_{H^\star}\longrightarrow0
\]
and
\[
  \operatorname{supp}v_m^+\subset[m,\infty)\times[0,d],
  \qquad
  \operatorname{supp}v_m^-\subset(-\infty,-m]\times[0,d].
\]
\end{lemma}
\begin{proof}
We give the right-end construction.  By the Weyl criterion for the spectrum,
choose $f_m\in D(A^+)$ with
\[
  \|f_m\|_{H^+}=1,\qquad
  \|(A^+-\lambda)f_m\|_{H^+}\leq m^{-1}.
\]
Since $A^+$ is nonnegative and associated with the gradient form,
\[
  \|\nabla f_m\|_{L^2(B)}^2
  =\langle A^+f_m,f_m\rangle_{H^+}
  \leq\lambda+m^{-1}
\]
for all sufficiently large $m$.  The periodic translation
\[
  (\tau_q^+f)(x,y):=f(x-qL^+,y),\qquad q\in\Z,
\]
is unitary in $H^+$ and commutes with $A^+$.  Choose $q_m$ so large that
\[
  \bigl\|\mathbf 1_{(-\infty,2m]}(Q)\tau_{q_m}^+f_m
  \bigr\|_{H^+}\leq m^{-1}.
\]
Let $\eta\in C^\infty(\R;[0,1])$ satisfy
$\eta(t)=0$ for $t\leq0$ and $\eta(t)=1$ for $t\geq1$, and set
$\eta_m^+(x):=\eta((x-m)/m)$.  Then
\[
  g_m:=\eta_m^+\tau_{q_m}^+f_m
\]
is supported in $[m,\infty)\times[0,d]$ and
$\|g_m\|_{H^+}\to1$.  The domain-preservation calculation in
\Cref{lem:appendix-local-compactness} gives $g_m\in D(A^+)$, and
\[
  (A^+-\lambda)g_m
  =\eta_m^+\tau_{q_m}^+(A^+-\lambda)f_m
   +[A^+,\eta_m^+]\tau_{q_m}^+f_m.
\]
The first term tends to zero.  Formula
\eqref{eq:appendix-commutator-formula}, the uniform $H^1$ bound, and
$\|\partial_x^r\eta_m^+\|_\infty=O(m^{-r})$ for $r=1,2$ show that the second
term also tends to zero.  Normalize $g_m$ to obtain $v_m^+$.  Translating by
negative integer multiples of $L^-$ and using
$\eta_m^-(x):=\eta((-x-m)/m)$ gives the left-end sequence.  This is the
second-order strip counterpart of the translation-and-cutoff construction in
the proof of
\cite[Proposition~3.10]{DeNittisEtAl2021}.
\end{proof}

\begin{proof}[Proof of \Cref{thm:two-ended-essential-spectrum}]
Floquet--Bloch decomposition applied separately to the two periodic
extensions gives, by \eqref{eq:background-bands},
\begin{equation}
  \sigma_{\mathrm{ess}}(A_0)
  =\sigma_{\mathrm{ess}}(A^-)
   \cup\sigma_{\mathrm{ess}}(A^+)
  =\sigma(A^-)\cup\sigma(A^+)
  =\Sigma_-(\beta)\cup\Sigma_+(\beta).
  \label{eq:appendix-free-essential-spectrum}
\end{equation}
Here the first equality is the elementary essential-spectrum formula for a
finite orthogonal direct sum, while the second is the periodic-spectrum result
of \cite[Proposition~3.1, especially equations~(3.3)--(3.4)]{Fliss2013}.

We first prove
$\sigma_{\mathrm{ess}}(A)\subset\sigma_{\mathrm{ess}}(A_0)$.
Choose $j_-,j_+\in C^\infty(\R;[0,1])$ such that
\[
\begin{array}{lll}
  j_-(x)=1&\text{for }x\leq-2R_0,&
  j_-(x)=0\text{ for }x\geq-R_0,\\
  j_+(x)=0&\text{for }x\leq R_0,&
  j_+(x)=1\text{ for }x\geq2R_0.
\end{array}
\]
Let $\lambda\in\sigma_{\mathrm{ess}}(A)$ and take the Zhislin sequence
$u_m$ furnished by \Cref{lem:appendix-Zhislin}.  For $m>2R_0$,
\[
  u_m=j_-u_m+j_+u_m,
\]
and the two terms have disjoint supports.  Define
\[
  u_m^0:=c_m(j_-u_m,j_+u_m)\in D(A_0),
  \qquad
  c_m:=\|(j_-u_m,j_+u_m)\|_{H_0}^{-1}.
\]
On the support of $j_-u_m$ the weights $\rho$ and $\rho^-$ agree, and on the
support of $j_+u_m$ the weights $\rho$ and $\rho^+$ agree.  Hence $c_m=1$ for
all sufficiently large $m$.  The derivatives of $j_\pm$ are supported in the
fixed set $[-2R_0,2R_0]$, where $u_m$ vanishes.  Therefore the commutator
terms vanish and
\[
\begin{aligned}
  \|(A_0-\lambda)u_m^0\|_{H_0}^2
  &=\|(A^--\lambda)j_-u_m\|_{H^-}^2
    +\|(A^+-\lambda)j_+u_m\|_{H^+}^2\\
  &\leq C\|(A-\lambda)u_m\|_H^2\longrightarrow0.
\end{aligned}
\]
The sequence $u_m^0$ still vanishes in $|x|\leq m$.  By
\Cref{lem:appendix-Zhislin},
$\lambda\in\sigma_{\mathrm{ess}}(A_0)$.

For the reverse inclusion, let
$\lambda\in\sigma(A^+)$.  By
\Cref{lem:appendix-one-sided-Weyl}, there is a normalized approximate
eigenvector sequence $v_m^+$ supported in
$[m,\infty)\times[0,d]$.  Once $m>R_0$, the two weights and the two
differential realizations agree on this support.  Thus
\[
  \|v_m^+\|_H=1,\qquad
  \|(A-\lambda)v_m^+\|_H
  =\|(A^+-\lambda)v_m^+\|_{H^+}\longrightarrow0.
\]
Because the supports escape to the right, $v_m^+\rightharpoonup0$ in $H$.
It is a singular Weyl sequence for $A$, and hence
$\lambda\in\sigma_{\mathrm{ess}}(A)$.  The left-end construction gives
$\sigma(A^-)\subset\sigma_{\mathrm{ess}}(A)$ in the same way.  Consequently,
\[
  \sigma(A^-)\cup\sigma(A^+)
  \subset\sigma_{\mathrm{ess}}(A)
  \subset\sigma_{\mathrm{ess}}(A_0).
\]
Combining this with \eqref{eq:appendix-free-essential-spectrum} proves
\eqref{eq:projected-spectrum}.
\end{proof}
